% ==============================================================================
% 04_criteria1_plos.tex — Criterion 1: Expected Learning Outcomes
% Score: 2
% ==============================================================================

\criterion{1}{ผลการเรียนรู้ที่คาดหวัง}{Expected Learning Outcomes}

% ------------------------------------------------------------------------------
\criterionsub{1.1 The programme to show that the expected learning outcomes are appropriately formulated in accordance with an established learning taxonomy, are aligned to the vision and mission of the university, and are known to all stakeholders.}

หลักสูตรวิศวกรรมศาสตรมหาบัณฑิต สาขาวิชาวิศวกรรมคอมพิวเตอร์และปัญญาประดิษฐ์ (หลักสูตรใหม่ พ.ศ.~2568) ของมหาวิทยาลัยราชภัฏอุตรดิตถ์ ได้ดำเนินการกำหนดและพัฒนาผลการเรียนรู้ระดับหลักสูตร (Programme Learning Outcomes: PLOs) ภายใต้กรอบมาตรฐานคุณภาพการศึกษาระดับบัณฑิตศึกษาและเกณฑ์การประเมินคุณภาพหลักสูตร โดยนำความเห็นของผู้มีส่วนได้ส่วนเสียเข้ามาประกอบการตัดสินใจในการยกร่าง และผ่านกลไกกำกับคุณภาพระดับคณะและมหาวิทยาลัยตามลำดับ ตั้งแต่การให้ความเห็นชอบในระดับคณะ การกลั่นกรองโดยคณะกรรมการบัณฑิตศึกษาประจำมหาวิทยาลัย การพิจารณาโดยสภาวิชาการ ก่อนเสนอสภามหาวิทยาลัยเพื่ออนุมัติ (มติครั้งที่ 4/2568 ลงวันที่ 4 เมษายน พ.ศ.~2568) เพื่อเตรียมการเปิดรับนักศึกษาในปีการศึกษา 2568 (ภาคการศึกษาที่ 2 มีนักศึกษาเข้าศึกษาจริง 4 คน) กรอบ PLO ที่ระบุในมคอ.2 จึงมีความชัดเจน สอดคล้องกับนโยบายมหาวิทยาลัย และตรวจสอบย้อนกลับได้จากเอกสารการประชุมที่เกี่ยวข้อง (อ้างอิง \hyperlink{ref:AUNQA-1-2}{AUNQA-1-2} และ \hyperlink{ref:AUNQA-1-4}{AUNQA-1-4})

การเก็บรวบรวมข้อมูลความเห็นของผู้มีส่วนได้ส่วนเสียครอบคลุม~5 กลุ่มหลักผ่านวิธีการที่หลากหลาย กลุ่มผู้ใช้บัณฑิตและสถานประกอบการด้านดิจิทัล/AI ในภาคเหนือตอนล่างให้ข้อมูลผ่านแบบสอบถามและการสนทนากลุ่มแบบมีโครงสร้างในช่วงยกร่างหลักสูตร กลุ่มศิษย์เก่าหลักสูตรวิศวกรรมคอมพิวเตอร์ระดับปริญญาตรีให้ข้อมูลผ่านการสนทนากลุ่มร่วมกับการทบทวนแนวโน้มอุตสาหกรรมจากเอกสารนโยบายกำลังคน AI ของภาครัฐ ส่วนองค์กรปกครองส่วนท้องถิ่นและชุมชนให้ข้อมูลผ่านการหารือและการมีส่วนร่วมของหน่วยงาน (อปท. และเครือข่ายพันธมิตร) สะท้อนความคาดหวังด้านการพัฒนากำลังคนท้องถิ่น สำหรับกลุ่มนักศึกษาและผู้สมัครเข้าศึกษา หลักสูตรเผยแพร่ข้อมูลและ PLO ผ่านเอกสารรับสมัครและแผนการปฐมนิเทศ ขณะที่อาจารย์ผู้รับผิดชอบหลักสูตรให้ข้อมูลผ่านการประชุมคณะกรรมการพัฒนาหลักสูตรและการทบทวนร่าง PLO/PEO อ้างอิงกรอบ AUN-QA v4.0 และมาตรฐานคุณวุฒิระดับบัณฑิตศึกษา (TQF) นอกจากนี้หลักสูตรได้เผยแพร่รายละเอียดหลักสูตรตาม มคอ.2 แก่กลุ่มเป้าหมายผ่านเว็บไซต์มหาวิทยาลัยและคณะเทคโนโลยีอุตสาหกรรม เพจ Facebook ของหลักสูตร ``\href{https://web.facebook.com/profile.php?id=61577236346855}{วิศวกรรมคอมพิวเตอร์และปัญญาประดิษฐ์}'' เอกสารแนะนำหลักสูตรสำหรับผู้สมัคร และช่องทางประชาสัมพันธ์ของมหาวิทยาลัย/คณะ

การกำหนดผลการเรียนรู้ของหลักสูตรจัดทำให้สอดคล้องกับกรอบคุณวุฒิแห่งชาติ (มคอ.) สำหรับระดับบัณฑิตศึกษา ครอบคลุมห้าด้าน ได้แก่ (1)~ด้านคุณธรรมจริยธรรม (2)~ด้านความรู้ (3)~ด้านทักษะทางปัญญา (4)~ด้านทักษะความสัมพันธ์และความรับผิดชอบ และ (5)~ด้านทักษะการวิเคราะห์เชิงตัวเลข การสื่อสาร และการใช้เทคโนโลยีสารสนเทศ ซึ่งข้อความ PLO กระจายการเรียนรู้ครบทุกด้านตามเจตนาของกรอบมคอ. จากนั้นหลักสูตรนำความต้องการที่สำคัญจากการรับฟังผู้มีส่วนได้ส่วนเสียไปวิเคราะห์และจัดลำดับความสำคัญ และเชื่อมโยงกับ \textbf{วิสัยทัศน์ พันธกิจ อัตลักษณ์ และปรัชญาการจัดการศึกษา} ของมหาวิทยาลัยและคณะที่ร่วมรับผิดชอบ จากนั้นกำหนดข้อความ PLO โดยอาศัยหลัก \textbf{Bloom's Revised Taxonomy} ควบคู่กับการเขียนวัตถุประสงค์ในรูปแบบ ABCD (Action verb, Behavior, Condition, Degree) เพื่อให้คำกริยาและผลลัพธ์มีความชัดเจนและวัดผลได้ สอดคล้องกับหลักปฏิบัติด้านประกันคุณภาพการศึกษาระหว่างประเทศและในประเทศ โดยเฉพาะอย่างยิ่งการจัดความสัมพันธ์ของ PLO ต่อวิสัยทัศน์ ``มหาวิทยาลัยพันธกิจสัมพันธ์ที่มีคุณภาพ สร้างคุณค่าเพื่อพัฒนาท้องถิ่น'' และปรัชญา IEL (Integrated Experiential Learning)

\begin{table}[H]
  \centering
  \caption{การสะท้อนเชิงธีมของ PLO ต่อวิสัยทัศน์ พันธกิจ/ภารกิจด้านผลผลิต และปรัชญา IEL}
  \label{tbl:plo_vision_alignment}
  \begingroup
  \small
  \setlength{\tabcolsep}{3pt}
  \renewcommand{\arraystretch}{1.2}
  \begin{tabularx}{\linewidth}{|>{\raggedright\arraybackslash}X|c|c|c|c|c|}
    \hline
    \textbf{ข้อความทิศทาง (สรุป)} & \textbf{PLO1} & \textbf{PLO2} & \textbf{PLO3} & \textbf{PLO4} & \textbf{PLO5} \\
    \hline
    วิสัยทัศน์ มร.อ.: คุณภาพ พันธกิจสัมพันธ์ พัฒนาท้องถิ่น & \checkmark & \checkmark & \checkmark & \checkmark & \checkmark \\
    \hline
    พันธกิจ/ความจำเป็น: กำลังคนดิจิทัล--ปัญญาประดิษฐ์ และการแก้ปัญหาในพื้นที่ & \checkmark & \checkmark & \checkmark &  & \checkmark \\
    \hline
    ปรัชญา IEL: ประสบการณ์เชิงบูรณาการ & \checkmark & \checkmark & \checkmark & \checkmark & \checkmark \\
    \hline
  \end{tabularx}
  \endgroup
\end{table}

\begin{table}[H]
  \centering
  \caption{Bloom's Level ของแต่ละ PLO}
  \label{tbl:bloom_plo}
  \begin{tabularx}{\linewidth}{|>{\raggedright\arraybackslash}X|>{\raggedright\arraybackslash}X|>{\raggedright\arraybackslash}X|>{\raggedright\arraybackslash}X|}
    \hline
    \textbf{PLO} & \textbf{ชื่อย่อ} & \textbf{Bloom's Level} & \textbf{ระดับความซับซ้อน} \\
    \hline
    PLO1 & Create           & Create (C6)   & สร้างสรรค์สูงสุด \\
    PLO2 & Apply            & Apply (C3)    & ประยุกต์ใช้ \\
    PLO3 & Create Knowledge & Create (C6)   & สร้างสรรค์สูงสุด \\
    PLO4 & Ethics           & Evaluate (C5) & ประเมินค่า \\
    PLO5 & Analyze/Comm.    & Analyze (C4)  & วิเคราะห์ \\
    \hline
  \end{tabularx}
\end{table}

\noindent\textbf{ข้อความผลการเรียนรู้ที่คาดหวัง (PLOs) ของหลักสูตร}

ทั้งสองแผนการศึกษากำหนด PLO ไว้ 5 ข้อเหมือนกัน โดย PLO1, PLO2, PLO4 และ PLO5 มีข้อความตรงกัน ส่วน PLO3 ต่างกันตามจุดเน้นของแต่ละแผน (มคอ.2 P267--P279)

\textbf{แผนที่ 1 (วิทยานิพนธ์) -- PLO1--PLO5}
\begin{enumerate}
  \item \textbf{PLO1 (Create)} ใช้ซอฟต์แวร์หรือฮาร์ดแวร์ทางคอมพิวเตอร์ในการสร้างสรรค์ผลงาน องค์ความรู้ และนวัตกรรม
  \item \textbf{PLO2 (Apply)} ประยุกต์ใช้ความรู้ในสาขาวิชาวิศวกรรมคอมพิวเตอร์เพื่อแก้ไขปัญหาความต้องการขององค์กร ชุมชน และท้องถิ่น
  \item \textbf{PLO3 (Create Knowledge)} สร้างองค์ความรู้ทางด้านวิศวกรรมคอมพิวเตอร์และปัญญาประดิษฐ์ที่ตอบสนองความต้องการขององค์กร
  \item \textbf{PLO4 (Ethics)} ปฏิบัติตามหลักจรรยาบรรณวิชาชีพวิศวกรรมคอมพิวเตอร์
  \item \textbf{PLO5 (Analyze/Communicate/Collaborate/Lifelong)} สามารถคิดวิเคราะห์ การสื่อสาร การทำงานร่วมกัน การปรับตัว และการเรียนรู้ตลอดชีวิต
\end{enumerate}

\textbf{แผนที่ 2 (สารนิพนธ์) -- PLO1--PLO5}
\begin{enumerate}
  \item \textbf{PLO1 (Create)} ใช้ซอฟต์แวร์หรือฮาร์ดแวร์ทางคอมพิวเตอร์ในการสร้างสรรค์ผลงาน องค์ความรู้ และนวัตกรรม
  \item \textbf{PLO2 (Apply)} ประยุกต์ใช้ความรู้ในสาขาวิชาวิศวกรรมคอมพิวเตอร์เพื่อแก้ไขปัญหาความต้องการขององค์กร ชุมชน และท้องถิ่น
  \item \textbf{PLO3 (Create Prototype)} สร้างต้นแบบนวัตกรรมด้านวิศวกรรมคอมพิวเตอร์และปัญญาประดิษฐ์ที่ตอบสนองความต้องการขององค์กร
  \item \textbf{PLO4 (Ethics)} ปฏิบัติตามหลักจรรยาบรรณวิชาชีพวิศวกรรมคอมพิวเตอร์
  \item \textbf{PLO5 (Analyze/Communicate/Collaborate/Lifelong)} สามารถคิดวิเคราะห์ การสื่อสาร การทำงานร่วมกัน การปรับตัว และการเรียนรู้ตลอดชีวิต
\end{enumerate}

PLO3 ของแผน 1 มุ่งเน้นการสร้างองค์ความรู้ใหม่ทางวิชาการ ส่วน PLO3 ของแผน 2 มุ่งเน้นการพัฒนาต้นแบบนวัตกรรมเชิงประยุกต์ การวัดผล PLO3 ในชั้นปีที่ 2 จะแยกเกณฑ์ตามรูบริกของวิทยานิพนธ์และสารนิพนธ์ตามลำดับ


\begin{table}[H]
  \centering
  \caption{ความเชื่อมโยง PEO--PLO (ทั้งสองแผนการศึกษา)}
  \label{tbl:peo_plo_both}
  \begingroup
  \small
  \setlength{\tabcolsep}{3pt}
  \renewcommand{\arraystretch}{1.2}
  \begin{tabularx}{\linewidth}{|>{\raggedright\arraybackslash}X|c|c|c|c|c|}
    \hline
    \rowcolor{gray!15}
    \textbf{PEO (ทั้ง 2 แผน)} & \textbf{PLO1} & \textbf{PLO2} & \textbf{PLO3} & \textbf{PLO4} & \textbf{PLO5} \\
    \hline
    PEO1: ทักษะซอฟต์แวร์/ฮาร์ดแวร์ในการออกแบบและพัฒนาองค์ความรู้ & \checkmark &            &            &            &            \\
    \hline
    PEO2 แผน 1: พัฒนาองค์ความรู้แก้ปัญหาในองค์กร (CE\&AI) &            & \checkmark & \checkmark &            &            \\
    \hline
    PEO2 แผน 2: พัฒนานวัตกรรมแก้ปัญหาในองค์กร (CE\&AI)   &            & \checkmark & \checkmark &            &            \\
    \hline
    PEO3: คุณธรรมและจรรยาบรรณวิชาชีพทางวิชาการ &            &            &            & \checkmark &            \\
    \hline
    PEO4: ทักษะการพัฒนาตนเองและเรียนรู้ตลอดชีวิต &            &            &            &            & \checkmark \\
    \hline
  \end{tabularx}
  \endgroup
\end{table}

\noindent PEO1, PEO3 และ PEO4 มีข้อความเหมือนกันทั้งสองแผน ส่วน PEO2 ต่างกันตามจุดเน้น: แผน 1 เน้นการพัฒนาองค์ความรู้ (วิจัย) ส่วนแผน 2 เน้นการพัฒนานวัตกรรม (ประยุกต์) ซึ่งสะท้อนออกสู่ PLO3 ที่ต่างกันระหว่างสองแผนเช่นกัน

% ------------------------------------------------------------------------------
\criterionsub{1.2 The programme to show that the expected learning outcomes for all courses are appropriately formulated and are aligned to the expected learning outcomes of the programme.}

ผลการเรียนรู้ของหลักสูตร (PLOs) ได้ถูกถ่ายทอดกระจายไปสู่รายวิชาบังคับ 6 รายวิชา และวิทยานิพนธ์/สารนิพนธ์ ผ่านกระบวนการทำ CLO (Course Learning Outcomes) โดยอาจารย์ผู้รับผิดชอบรายวิชาแต่ละท่าน กำหนด CLO ตามหลัก ABCD และเชื่อมโยง CLO เข้ากับ PLO อย่างน้อย 1 ข้อ ผ่านตาราง CLO--PLO Mapping ที่บรรจุไว้ใน มคอ.2 และในเอกสาร ``รวมเล่ม CLO รายวิชา'' (\hyperlink{ref:AUNQA-1-3}{AUNQA-1-3})

\begin{table}[H]
  \centering
  \caption{การกระจายการเรียนรู้ที่คาดหวังของหลักสูตรลงสู่รายวิชาบังคับ}
  \begingroup
  \fontbody
  \setlength{\tabcolsep}{5pt}
  \renewcommand{\arraystretch}{1.2}
  \begin{tabularx}{\linewidth}{|>{\raggedright\arraybackslash}X|c|c|c|c|c|}
    \hline
    \textbf{รายวิชา} & \textbf{PLO1} & \textbf{PLO2} & \textbf{PLO3} & \textbf{PLO4} & \textbf{PLO5} \\
    \hline
    4095101 ปัญญาประดิษฐ์และการเรียนรู้ของเครื่อง & \checkmark &            & \checkmark &            &            \\
    4095102 การเขียนโปรแกรมขั้นสูงสำหรับ ML       & \checkmark & \checkmark & \checkmark & \checkmark &            \\
    7015906 ระเบียบวิธีวิจัยฯ                      &            & \checkmark & \checkmark & \checkmark & \checkmark \\
    7015907 สัมมนาวิศวกรรม CPE-AI                   &            & \checkmark & \checkmark & \checkmark & \checkmark \\
    7015101 เทคโนโลยีอุบัติใหม่ฯ                   & \checkmark & \checkmark &            &            &            \\
    7015908 หัวข้อพิเศษฯ                           & \checkmark & \checkmark &            & \checkmark &            \\
    \hline
  \end{tabularx}
  \endgroup
\end{table}

\noindent\textbf{ข้อสังเกตเกี่ยวกับระดับพัฒนาการ (Learning Progression):} เป็นที่น่าสังเกตว่าในระดับรายวิชาบังคับชั้นปีที่ 1 (เช่น วิชา 4095101 และ 4095102) กำหนดระดับพุทธิพิสัยสูงสุดของ CLO อยู่ที่ระดับวิเคราะห์ (Analyze: L4) หรือต่ำกว่า ซึ่งสะท้อนการสร้างพื้นฐานความรู้และทักษะที่จำเป็น อย่างไรก็ตาม ระดับผลการเรียนรู้ที่คาดหวังระดับหลักสูตร PLO1 และ PLO3 กำหนด Bloom's Level ไว้สูงถึงระดับสร้างสรรค์ (Create: L6) ซึ่งการบรรลุ L6 นี้จะเกิดขึ้นอย่างสมบูรณ์และวัดผลได้จริงผ่านการดำเนินงานวิทยานิพนธ์หรือสารนิพนธ์ในชั้นปีที่ 2 (เช่น รายวิชา 7015901--7015903 และ 7015904--7015905) ตามกรอบการจัดระบบการเรียนรู้แบบไต่ระดับ (OBE Spiral Progression)


% ------------------------------------------------------------------------------
\criterionsub{1.3 The programme to show that the expected learning outcomes consist of both generic outcomes (related to written and oral communication, problem-solving, information technology, teambuilding skills, etc) and subject specific outcomes (related to knowledge and skills of the study discipline).}

ผลการเรียนรู้ที่คาดหวังของหลักสูตรประกอบด้วย Generic Learning Outcomes และ Subject Specific Learning Outcomes ดังตาราง

\begin{table}[H]
  \centering
  \caption{Generic Learning Outcomes และ Subject Specific Learning Outcomes}
  \begin{tabularx}{\linewidth}{|c|Y|c|c|}
    \hline
    \rowcolor{gray!15}
    \textbf{PLO} & \textbf{คำอธิบายผลการเรียนรู้} & \textbf{ประเภท} & \textbf{Level} \\
    \hline
    PLO1 & ใช้ซอฟต์แวร์/ฮาร์ดแวร์สร้างสรรค์ผลงาน นวัตกรรม & Subject specific & C \\
    \hline
    PLO2 & ประยุกต์ความรู้แก้ปัญหาขององค์กร/ชุมชน/ท้องถิ่น & Subject specific & A \\
    \hline
    PLO3 & สร้างองค์ความรู้/ต้นแบบนวัตกรรม CE \& AI & Subject specific & C \\
    \hline
    PLO4 & ปฏิบัติตามจรรยาบรรณวิชาชีพวิศวกรรมคอมพิวเตอร์ & Generic & E \\
    \hline
    PLO5 & คิดวิเคราะห์ สื่อสาร ทำงานร่วมกัน เรียนรู้ตลอดชีวิต & Generic & AN \\
    \hline
  \end{tabularx}
\end{table}

\noindent\textbf{หมายเหตุ Level learning:} R = Remembering, U = Understanding, A = Applying, AN = Analyzing, E = Evaluating, C = Creating

% ------------------------------------------------------------------------------
\criterionsub{1.4 The programme to show that the requirements of the stakeholders, especially the external stakeholders, are gathered, and that these are reflected in the expected learning outcomes.}

หลักสูตรกำหนดผู้มีส่วนได้ส่วนเสียออกเป็น 4 กลุ่มหลัก โดยสามารถจัดเป็นสองมิติ ได้แก่ \textbf{ผู้มีส่วนได้ส่วนเสียภายใน} (อาจารย์ผู้รับผิดชอบหลักสูตร/อาจารย์ผู้สอน และนักศึกษาเมื่อเริ่มเปิดรับแล้ว) และ \textbf{ผู้มีส่วนได้ส่วนเสียภายนอก} (ผู้ใช้บัณฑิตในสถานประกอบการด้านดิจิทัล/AI ในภาคเหนือตอนล่าง องค์กรปกครองส่วนท้องถิ่นในจังหวัดอุตรดิตถ์และจังหวัดใกล้เคียง รวมถึงศิษย์เก่าหลักสูตรวิศวกรรมคอมพิวเตอร์ระดับปริญญาตรี) เพื่อให้การรวบรวมความต้องการครอบคลุมทั้งมุมมองผู้จัดการเรียนรู้ ผู้เรียน และผู้ใช้บัณฑิตในระบบเศรษฐกิจและสังคม

วิธีการได้มาซึ่งความต้องการของผู้มีส่วนได้ส่วนเสียประกอบด้วย (1) แบบสอบถามออนไลน์ผู้มีส่วนได้ส่วนเสียภายนอก (2) การสนทนากลุ่ม (focus group) กับศิษย์เก่าและผู้ใช้บัณฑิต (3) การวิเคราะห์แนวโน้มอุตสาหกรรมจากเอกสาร AI Workforce Development Plan ของรัฐบาล หลังจากนั้นนำข้อมูลมาวิเคราะห์/จัดกลุ่มความต้องการ และสะท้อนเข้าสู่ PLOs

\textbf{ผลการสำรวจความต้องการของผู้มีส่วนได้ส่วนเสีย (\hyperlink{ref:AUNQA-1-4b}{AUNQA-1-4b})} ดำเนินการเก็บข้อมูลผ่านแบบสอบถามออนไลน์ในช่วงเดือนมิถุนายน~2567 ได้รับแบบสอบถามคืน $n = 31$ ราย ประกอบด้วยผู้ที่กำลังศึกษาอยู่ 17~คน (54.8\%) และผู้สำเร็จการศึกษาและทำงานแล้ว 14~คน (45.2\%) มีสัดส่วนเพศชาย 21~คน (67.7\%) และเพศหญิง 10~คน (32.3\%) กระจายช่วงอายุตั้งแต่ 18 ปีจนถึงมากกว่า 35~ปี

\begin{table}[H]
  \centering
  \caption{สรุปผลการสำรวจความต้องการผู้มีส่วนได้ส่วนเสีย ($n = 31$, มิถุนายน 2567)}
  \label{tbl:stakeholder_survey}
  \begingroup
  \small
  \setlength{\tabcolsep}{4pt}
  \renewcommand{\arraystretch}{1.2}
  \begin{tabularx}{\linewidth}{|>{\raggedright\arraybackslash}X|r|r|}
    \hline
    \rowcolor{gray!15}
    \textbf{ประเด็น / ผลการสำรวจ} & \textbf{จำนวน (คน)} & \textbf{ร้อยละ} \\
    \hline
    \multicolumn{3}{|l|}{\textit{ระดับการศึกษาที่ต้องการ}} \\
    \hline
    ระดับปริญญาโท (ป.โท 2 ปี) & 14 & 45.2 \\
    ระดับปริญาตรีต่อเนื่องปริญญาโท & 5  & 16.1 \\
    ระดับปริญญาเอก & 6  & 19.4 \\
    ไม่ต้องการศึกษาต่อ & 6  & 19.3 \\
    \hline
    \multicolumn{3}{|l|}{\textit{เป้าหมายในการศึกษาต่อ (เลือกได้มากกว่า 1 ข้อ)}} \\
    \hline
    ต้องการต่อยอดองค์ความรู้ในสาขาที่สนใจ & 14 & 45.2 \\
    เพื่อใช้ในการประกอบอาชีพในอนาคต & 13 & 41.9 \\
    แก้ปัญหาหรือพัฒนางานที่ทำอยู่ & 7  & 22.6 \\
    ปรับเพิ่มวุฒิการศึกษา / เปลี่ยนตำแหน่งงาน & 7  & 22.6 \\
    \hline
    \multicolumn{3}{|l|}{\textit{ช่วงเวลาเรียนที่เหมาะสม}} \\
    \hline
    วันเสาร์-อาทิตย์ (เฉพาะสุดสัปดาห์) & 26 & 83.9 \\
    วันจันทร์-ศุกร์ (วันธรรมดา) & 5  & 16.1 \\
    \hline
    \multicolumn{3}{|l|}{\textit{แผนการศึกษาต่อ}} \\
    \hline
    ต้องการศึกษาต่อทันที & 11 & 35.5 \\
    มีแผนอีก 1--2 ปีข้างหน้า & 8  & 25.8 \\
    มีแผนอีก 3 ปีข้างหน้า & 5  & 16.1 \\
    ไม่มีแผนศึกษาต่อ & 7  & 22.6 \\
    \hline
  \end{tabularx}
  \endgroup
\end{table}

\begin{table}[H]
  \centering
  \caption{การสะท้อนความต้องการผู้มีส่วนได้ส่วนเสียลงสู่ PLO}
  \label{tbl:sh_to_plo}
  \begingroup
  \small
  \setlength{\tabcolsep}{4pt}
  \renewcommand{\arraystretch}{1.2}
  \begin{tabularx}{\linewidth}{|>{\raggedright\arraybackslash}X|c|c|c|c|c|}
    \hline
    \rowcolor{gray!15}
    \textbf{ความต้องการของผู้มีส่วนได้ส่วนเสีย} & \textbf{PLO1} & \textbf{PLO2} & \textbf{PLO3} & \textbf{PLO4} & \textbf{PLO5} \\
    \hline
    ทักษะสร้างซอฟต์แวร์/ระบบ AI และนวัตกรรมดิจิทัล & \checkmark & & \checkmark & & \\
    \hline
    ประยุกต์ความรู้แก้ปัญหาในองค์กร/ท้องถิ่น & & \checkmark & & & \\
    \hline
    สร้างองค์ความรู้ใหม่/ต่อยอดการวิจัย & & & \checkmark & & \\
    \hline
    จริยธรรมวิชาชีพ/มาตรฐานวิศวกร & & & & \checkmark & \\
    \hline
    ทักษะการสื่อสาร ทำงานร่วมกัน และเรียนรู้ตลอดชีวิต & & & & & \checkmark \\
    \hline
  \end{tabularx}
  \endgroup
\end{table}

ผลการสำรวจยืนยันว่ากลุ่มเป้าหมายหลักของหลักสูตรต้องการระดับปริญญาโท รวม 61.3\% (19~คน) และมีความต้องการเรียนในรูปแบบ Weekend class (เสาร์-อาทิตย์) ถึง 83.9\% ซึ่งสอดคล้องกับรูปแบบการจัดการเรียนการสอนที่หลักสูตรวางแผนไว้ เป้าหมายหลักในการศึกษาต่อ ได้แก่ การต่อยอดองค์ความรู้ (45.2\%) การพัฒนาอาชีพ (41.9\%) และการแก้ปัญหาในงาน (22.6\%) สะท้อนตรงกับ PLO1--PLO3 ที่เน้นการสร้างนวัตกรรม การแก้ปัญหา และการสร้างองค์ความรู้ใหม่

% ------------------------------------------------------------------------------
\criterionsub{1.5 The programme to show that the expected learning outcomes are achieved by the students by the time they graduate.}

หลักสูตรเป็นหลักสูตรใหม่ พ.ศ. 2568 เริ่มมีนักศึกษาเข้าศึกษา (ปีการศึกษา 2568 จำนวน 4 คน ภาคการศึกษาที่ 2) จึงยังไม่มีผลการบรรลุ PLOs ของผู้สำเร็จการศึกษาระดับหลักสูตร อย่างไรก็ตาม หลักสูตรได้วางระบบการวัดการบรรลุ PLOs ครบถ้วนตามกรอบ 4 มิติ (When / Method / Tool / Person) ดังตาราง~\ref{tbl:plo_achievement_plan}

\begin{table}[H]
  \centering
  \caption{แผนการวัดการบรรลุ PLOs ของหลักสูตร (PLO Achievement Measurement Plan)}
  \label{tbl:plo_achievement_plan}
  \begingroup
  \footnotesize
  \setlength{\tabcolsep}{3pt}
  \renewcommand{\arraystretch}{1.2}
  \begin{tabularx}{\linewidth}{|c|>{\raggedright\arraybackslash}X|>{\raggedright\arraybackslash}X|>{\raggedright\arraybackslash}X|>{\raggedright\arraybackslash}X|}
    \hline
    \textbf{PLO} & \textbf{When} & \textbf{Method} & \textbf{Tool} & \textbf{Person} \\
    \hline
    PLO1 & ปีการศึกษาสุดท้าย (สอบวิทยานิพนธ์) + CLO assessment รายวิชา & CLO aggregation + Thesis Defense & Rubric CLO (4 ระดับ) + เกณฑ์ผ่าน Thesis & อาจารย์ผู้สอน + คณะกรรมการสอบ \\
    \hline
    PLO2 & ปีการศึกษาที่ 1 ภาค~1--2 (วิชา 4095102, 7015101, 7015906) + ปีสุดท้าย & CLO assessment + Community Project & Rubric CLO + ผลประเมินโครงงานชุมชน & อาจารย์ผู้สอน \\
    \hline
    PLO3 & ปีการศึกษาสุดท้าย (สอบวิทยานิพนธ์/สารนิพนธ์) & Thesis/Dissertation Defense & เกณฑ์ผ่านการสอบปากเปล่า (External + Internal examiner) & คณะกรรมการสอบ (Internal + External) \\
    \hline
    PLO4 & ปีการศึกษาที่ 1 ภาค~1--2 (วิชา 4095102, 7015906) + ก่อนสำเร็จการศึกษา & CLO4 assessment + Exit survey จริยธรรม & Rubric Affective (CLO4) + แบบประเมินจริยธรรมวิจัย & อาจารย์ผู้สอน + อาจารย์ที่ปรึกษา \\
    \hline
    PLO5 & ปีการศึกษาที่ 1 ภาค~2 (วิชา 7015906) + ก่อนสำเร็จการศึกษา & CLO5 assessment + Self-assessment & Rubric CLO5 + แบบประเมิน LLL ตนเอง & อาจารย์ผู้สอน + Self-assessment \\
    \hline
  \end{tabularx}
  \endgroup
\end{table}

แม้ยังไม่มีบัณฑิตสำเร็จการศึกษา หลักสูตรได้ดำเนินการรวบรวมข้อมูล CLO Achievement จาก มคอ.5 ทั้ง 4 วิชาบังคับและ Aggregate ตาม CLO--PLO Mapping แล้ว ผลในปีการศึกษา~2568 ปรากฏว่า PLO ทั้ง 5 ตัวผ่านเกณฑ์หลักสูตร ($\geq$80\% ของนักศึกษา) ครบทุกตัว โดย PLO1--PLO5 ผ่านในอัตรา 100\% (4/4 คน) (อ้างอิง~\hyperlink{ref:AUNQA-8-4-1}{AUNQA-8-4-1}) ซึ่งยืนยันว่าระบบวัด PLO ที่ออกแบบไว้ทำงานได้ตามวัตถุประสงค์ในปีแรก ข้อมูล PLO Achievement ระดับบัณฑิต (Graduation-level) จะรายงานได้เมื่อนักศึกษารุ่นแรกสำเร็จการศึกษา ซึ่งคาดการณ์ไว้ในปีการศึกษา 2569--2570 (คาดภาค~1/2570 หรือภาค~2/2570 ขึ้นอยู่กับความก้าวหน้าวิทยานิพนธ์รายบุคคล)

\begin{strengthbox}
หลักสูตรกำหนด PLOs ครบถ้วนตาม Bloom's Revised Taxonomy (C3--C6) ครอบคลุมทั้ง Subject-specific (PLO1--3) และ Generic outcomes (PLO4--5) เชื่อมโยงกับวิสัยทัศน์/พันธกิจมหาวิทยาลัยและปรัชญา IEL อย่างเป็นระบบ และกระจาย PLO ลงสู่ CLO ทุกรายวิชาผ่านตาราง CLO--PLO Mapping ที่ชัดเจน นอกจากนี้ผลสำรวจความพึงพอใจนักศึกษาและผู้มีส่วนได้ส่วนเสีย ($n=4$, ปีการศึกษา~2568) ให้คะแนนด้านวัตถุประสงค์/PLO เฉลี่ย~4.75/5.00 สะท้อนว่า PLO มีความชัดเจนและสอดคล้องกับความคาดหวังของผู้เรียนจริง
\end{strengthbox}

\begin{weaknessbox}
ผลการบรรลุ PLO ในปีการศึกษา 2568 วัดผ่าน CLO Aggregation จาก 4 วิชาบังคับเท่านั้น ยังไม่ครอบคลุม PLO3 ในมิติการทำวิทยานิพนธ์/สารนิพนธ์จริง (รอการสอบป้องกัน ปีการศึกษา 2569--2570) นอกจากนี้ผลสำรวจ stakeholder ($n=31$) ก่อนเปิดหลักสูตรเป็นกลุ่มผู้สนใจทั่วไปและนักศึกษา ป.ตรี ยังไม่ใช่ผู้ใช้บัณฑิตระดับ ป.โท โดยตรง --- Formal Employer Survey จะดำเนินการเมื่อบัณฑิตรุ่นแรกสำเร็จการศึกษา (คาดการณ์ปีการศึกษา 2569--2570)
\end{weaknessbox}

\begin{selfassessmentbox}{2}{หลักสูตรกำหนด PLOs ครบถ้วนตาม Bloom's Taxonomy (C3--C6) เชื่อมโยงกับวิสัยทัศน์/พันธกิจ/IEL กระจายลงสู่รายวิชาผ่าน CLO--PLO Mapping มีหลักฐานเชิงปริมาณจากการสำรวจผู้มีส่วนได้ส่วนเสีย $n=31$ ราย ผลสำรวจความพึงพอใจนักศึกษา $n=4$ ปีการศึกษา~2568 คะแนนเฉลี่ย~4.75/5.00 และผลการวัด PLO Achievement รอบแรกผ่าน CLO Aggregation: PLO ทั้ง 5 ตัว 100\% (4/4 คน) มีแผนวัดการบรรลุ PLO ครบ 5 PLO ตามกรอบ When/Method/Tool/Person ข้อมูล PLO ระดับบัณฑิตจะรายงานในปีการศึกษา 2569--2570}
\end{selfassessmentbox}

\begin{evidencelist}
\evidence{AUNQA-1-1}{เล่ม มคอ.2 หลักสูตรวิศวกรรมคอมพิวเตอร์และปัญญาประดิษฐ์ ฉบับผ่านสภา\ \href{https://syweb.kidcbc.work/Eviden_aunqa68/\%E0\%B8\%A1\%E0\%B8\%84\%E0\%B8\%AD\%202\%20\%E0\%B8\%AB\%E0\%B8\%A5\%E0\%B8\%B1\%E0\%B8\%81\%E0\%B8\%AA\%E0\%B8\%B9\%E0\%B8\%95\%E0\%B8\%A3\%E0\%B8\%A7\%E0\%B8\%B4\%E0\%B8\%A8\%E0\%B8\%A7\%E0\%B8\%81\%E0\%B8\%A3\%E0\%B8\%A3\%E0\%B8\%A1\%E0\%B8\%84\%E0\%B8\%AD\%E0\%B8\%A1\%E0\%B8\%9E\%E0\%B8\%B4\%E0\%B8\%A7\%E0\%B9\%80\%E0\%B8\%95\%E0\%B8\%AD\%E0\%B8\%A3\%E0\%B9\%8C\%E0\%B9\%81\%E0\%B8\%A5\%E0\%B8\%B0\%E0\%B8\%9B\%E0\%B8\%B1\%E0\%B8\%8D\%E0\%B8\%8D\%E0\%B8\%B2\%E0\%B8\%9B\%E0\%B8\%A3\%E0\%B8\%B0\%E0\%B8\%94\%E0\%B8\%B4\%E0\%B8\%A9\%E0\%B8\%90\%E0\%B9\%8C\%20\%E0\%B8\%9C\%E0\%B9\%88\%E0\%B8\%B2\%E0\%B8\%99\%E0\%B8\%AA\%E0\%B8\%A0\%E0\%B8\%B2\%202\%20\%E0\%B9\%80\%E0\%B8\%9E\%E0\%B8\%B4\%E0\%B9\%88\%E0\%B8\%A1\%20Comment\%20\%E0\%B8\%AA\%E0\%B8\%A0\%E0\%B8\%B2\%20\%E0\%B9\%81\%E0\%B8\%A5\%E0\%B8\%B0\%20\%E0\%B8\%A1\%E0\%B8\%95\%E0\%B8\%B4\%E0\%B8\%AA\%E0\%B8\%A0\%E0\%B8\%B2.pdf}{\texttt{[PDF]}}}
\evidence{AUNQA-1-2}{รายงานการประชุมยกร่างหลักสูตร ครั้งที่ 1/2566 (15 ก.ย. 2566)\ \href{https://syweb.kidcbc.work/Eviden_aunqa68/\%E0\%B8\%AB\%E0\%B8\%A5\%E0\%B8\%B1\%E0\%B8\%81\%E0\%B8\%90\%E0\%B8\%B2\%E0\%B8\%99\%E0\%B8\%81\%E0\%B8\%B2\%E0\%B8\%A3\%E0\%B8\%9B\%E0\%B8\%A3\%E0\%B8\%B0\%E0\%B8\%8A\%E0\%B8\%B8\%E0\%B8\%A1/01_\%E0\%B8\%A3\%E0\%B8\%B0\%E0\%B8\%94\%E0\%B8\%B1\%E0\%B8\%9A\%E0\%B8\%AB\%E0\%B8\%A5\%E0\%B8\%B1\%E0\%B8\%81\%E0\%B8\%AA\%E0\%B8\%B9\%E0\%B8\%95\%E0\%B8\%A3/\%E0\%B8\%81\%E0\%B8\%A1\%E0\%B8\%A8_2566-09-15_\%E0\%B8\%A2\%E0\%B8\%81\%E0\%B8\%A3\%E0\%B9\%88\%E0\%B8\%B2\%E0\%B8\%87\%E0\%B8\%AB\%E0\%B8\%A5\%E0\%B8\%B1\%E0\%B8\%81\%E0\%B8\%AA\%E0\%B8\%B9\%E0\%B8\%95\%E0\%B8\%A3_\%E0\%B8\%84\%E0\%B8\%A3\%E0\%B8\%B1\%E0\%B9\%89\%E0\%B8\%87\%E0\%B8\%97\%E0\%B8\%B5\%E0\%B9\%881.pdf}{\texttt{[PDF]}}}
\evidence{AUNQA-1-3}{รวมเล่ม CLO รายวิชา พร้อม CLO--PLO Mapping\ \href{https://syweb.kidcbc.work/Eviden_aunqa68/\%E0\%B8\%A3\%E0\%B8\%A7\%E0\%B8\%A1\%E0\%B9\%80\%E0\%B8\%A5\%E0\%B9\%88\%E0\%B8\%A1\%20CLO\%20\%E0\%B8\%A3\%E0\%B8\%B2\%E0\%B8\%A2\%E0\%B8\%A7\%E0\%B8\%B4\%E0\%B8\%8A\%E0\%B8\%B2.pdf}{\texttt{[PDF]}}}
\evidence{AUNQA-1-4}{มติสภามหาวิทยาลัยราชภัฏอุตรดิตถ์ ครั้งที่ 4/2568\ \href{https://syweb.kidcbc.work/Eviden_aunqa68/\%E0\%B8\%AB\%E0\%B8\%A5\%E0\%B8\%B1\%E0\%B8\%81\%E0\%B8\%90\%E0\%B8\%B2\%E0\%B8\%99\%E0\%B8\%81\%E0\%B8\%B2\%E0\%B8\%A3\%E0\%B8\%9B\%E0\%B8\%A3\%E0\%B8\%B0\%E0\%B8\%8A\%E0\%B8\%B8\%E0\%B8\%A1/03_\%E0\%B8\%A3\%E0\%B8\%B0\%E0\%B8\%94\%E0\%B8\%B1\%E0\%B8\%9A\%E0\%B8\%A1\%E0\%B8\%AB\%E0\%B8\%B2\%E0\%B8\%A7\%E0\%B8\%B4\%E0\%B8\%97\%E0\%B8\%A2\%E0\%B8\%B2\%E0\%B8\%A5\%E0\%B8\%B1\%E0\%B8\%A2/\%E0\%B8\%A1\%E0\%B8\%A3\%E0\%B8\%AD_\%E0\%B8\%AA\%E0\%B8\%A0\%E0\%B8\%B2\%E0\%B8\%A1\%E0\%B8\%AB\%E0\%B8\%B2\%E0\%B8\%A7\%E0\%B8\%B4\%E0\%B8\%97\%E0\%B8\%A2\%E0\%B8\%B2\%E0\%B8\%A5\%E0\%B8\%B1\%E0\%B8\%A2_4-2568_\%E0\%B8\%AA\%E0\%B9\%81\%E0\%B8\%81\%E0\%B8\%99.pdf}{\texttt{[PDF]}}}
\evidence{AUNQA-1-4b}{แบบสอบถามผู้มีส่วนได้ส่วนเสีย: ผลสำรวจ $n=31$ ราย (มิถุนายน 2567)\ \href{https://syweb.kidcbc.work/Eviden_aunqa68/\%E0\%B8\%A3\%E0\%B8\%B0\%E0\%B9\%80\%E0\%B8\%9A\%E0\%B8\%B5\%E0\%B8\%A2\%E0\%B8\%9A\%E0\%B8\%82\%E0\%B9\%89\%E0\%B8\%AD\%E0\%B8\%9A\%E0\%B8\%B1\%E0\%B8\%87\%E0\%B8\%84\%E0\%B8\%B1\%E0\%B8\%9A\%E0\%B9\%81\%E0\%B8\%A5\%E0\%B8\%B0\%E0\%B8\%9B\%E0\%B8\%A3\%E0\%B8\%B0\%E0\%B8\%81\%E0\%B8\%B2\%E0\%B8\%A8\%E0\%B8\%AD\%E0\%B8\%B7\%E0\%B9\%88\%E0\%B8\%99\%E0\%B9\%86/AUNQA-1-4b_stakeholder_survey_responses_31\%E0\%B8\%84\%E0\%B8\%99.csv}{\texttt{[CSV]}}\ \href{https://docs.google.com/forms/d/1hcnZnYHP7y5AJxPN5nblqyS900z9UYTZ5Fs-E0JL-JQ/viewanalytics}{\texttt{[Google Report]}}}
\evidence{AUNQA-C6-STSAT-2568}{แบบสำรวจความพึงพอใจนักศึกษาและผู้มีส่วนได้ส่วนเสีย CPE\&AI ปีการศึกษา~2568 ($n=4$): คะแนนความพึงพอใจต่อวัตถุประสงค์/PLO เฉลี่ย~4.75/5.00\ \href{https://docs.google.com/forms/d/1xpfb6AfYkbOkidJNh6ci95ux4ZLYzjvpUuqARYjhr5k/viewanalytics}{\texttt{[Google Report]}}}
\evidence{AUNQA-8-4-1}{รายงานผลการประเมิน CLO--PLO Achievement ปีการศึกษา 2568 (PLO ทั้ง 5 ตัว 100\%, $n=4$ คน, 4 วิชาบังคับ)\ \href{https://syweb.kidcbc.work/Eviden_aunqa68/\%E0\%B8\%A3\%E0\%B8\%B0\%E0\%B9\%80\%E0\%B8\%9A\%E0\%B8\%B5\%E0\%B8\%A2\%E0\%B8\%9A\%E0\%B8\%82\%E0\%B9\%89\%E0\%B8\%AD\%E0\%B8\%9A\%E0\%B8\%B1\%E0\%B8\%87\%E0\%B8\%84\%E0\%B8\%B1\%E0\%B8\%9A\%E0\%B9\%81\%E0\%B8\%A5\%E0\%B8\%B0\%E0\%B8\%9B\%E0\%B8\%A3\%E0\%B8\%B0\%E0\%B8\%81\%E0\%B8\%B2\%E0\%B8\%A8\%E0\%B8\%AD\%E0\%B8\%B7\%E0\%B9\%88\%E0\%B8\%99\%E0\%B9\%86/AUNQA-8-4-1_CLO_PLO_Achievement_2568.pdf}{\texttt{[PDF]}}}
\end{evidencelist}
